% BHU PYQ -- Manually transcribed & error-corrected
% Paper: GRA/GRB-608: Regional Planning
% Exam: B.A./B.Sc. (Hons.) Semester VI Examination, 2022-23

\documentclass[11pt,a4paper]{article}
\usepackage[a4paper,top=2.5cm,bottom=2.5cm,left=2.8cm,right=2.8cm,headheight=14pt]{geometry}
\usepackage{amsmath,amssymb}
\usepackage[T1]{fontenc}
\usepackage[utf8]{inputenc}
\usepackage{lmodern,microtype,fancyhdr,enumitem,hyperref,lastpage,parskip}

\hypersetup{colorlinks=true,urlcolor=black,linkcolor=black,
  pdftitle={GRA/GRB-608 Regional Planning -- BHU B.A./B.Sc. Sem VI 2022-23}}
\pagestyle{fancy}\fancyhf{}
\renewcommand{\headrulewidth}{0.4pt}\renewcommand{\footrulewidth}{0.4pt}
\fancyhead[L]{\small\textit{Banaras Hindu University}}
\fancyhead[R]{\small\textit{GRA/GRB-608: Regional Planning}}
\fancyfoot[C]{\small Page \thepage\ of \pageref{LastPage}}

\newlist{parts}{enumerate}{2}
\setlist[parts,1]{label=(\alph*),leftmargin=*,itemsep=5pt,topsep=3pt}
\newcommand{\pts}[1]{\hfill\textit{[#1]}}

\begin{document}
\begin{center}
  {\large\bfseries BANARAS HINDU UNIVERSITY}\\[2pt]
  {\normalsize B.A./B.Sc. (Hons.) --- Semester VI Examination, 2022-23}\\[6pt]
  \rule{0.8\linewidth}{0.5pt}\\[6pt]
  {\large\bfseries Paper No. GRA/GRB-608: Regional Planning}\\[4pt]
  {\normalsize Time: 3 Hours \hfill Maximum Marks: 70}\\[2pt]
  {\small Ref. No.: 59449}
\end{center}
\rule{\linewidth}{0.4pt}
\smallskip
\noindent\textit{Instructions: Answer \textbf{five} questions in all. Question No. 1, carrying 20 marks, is compulsory. Remaining questions are of equal marks.}
\bigskip

\noindent\textbf{Question 1.} Write short note of the following: \pts{20}
\begin{parts}
  \item Scope of regional planning
  \item Role of participatory approach in regional planning
  \item Concept of growth pole
  \item Concept of regional development
  \item Evolution of regional planning in India
\end{parts}

\medskip\noindent\textbf{Question 2.} Discuss the different approaches of regional planning. \pts{12.5}

\medskip\noindent\textbf{Question 3.} Explain the concept of multi-level planning with special reference to India. \pts{12.5}

\medskip\noindent\textbf{Question 4.} Describe the causes and pattern of regional disparities in India. \pts{12.5}

\medskip\noindent\textbf{Question 5.} Discuss the role of industry and transport in regional development. \pts{12.5}

\medskip\noindent\textbf{Question 6.} Define a planning region and examine the delimitation of macro and meso planning regions. \pts{12.5}

\medskip\noindent\textbf{Question 7.} Discuss the developmental issues and planning strategies for Eastern Uttar Pradesh. \pts{12.5}

\medskip\noindent\textbf{Question 8.} Examine the merits and demerits of Panchayati Raj planning in India. \pts{12.5}

\vfill\rule{\linewidth}{0.4pt}
\begin{center}\small
\href{https://bhu-exam.vercel.app/}{Click here to visit our website}\\[4pt]
\href{https://me-aryan.vercel.app/}{Click here to visit our contributor}\\[4pt]
\href{https://quashan-me.vercel.app/}{Click here to visit our contributor}
\end{center}
\end{document}
